% GENERATED FILE. Edit canonical lessons, then run npm run generate:books.
% canonical-source-hash: fnv1a64:95a728663fbb6516
% canonical-lessons: GE-C98-bibliothek, GE-C98-fahrrad, GE-C98-haltestelle, GE-C98-kleidung, GE-C98-hemd

\chapter{Library, Bicycle, Bus or tram stop, Clothes, Shirt}
\label{ch:ge-library-bicycle-bus-or-tram-stop-clothes-shirt}
\hlchaptermodality{98}

\begin{chapteropening}
\textbf{By the end of this chapter:} \emph{I can say a library, a bicycle, a bus or tram stop, clothes and a shirt.}

Complete the last lesson of chapter 98: I can say a library, a bicycle, a bus or tram stop, clothes and a shirt.
\end{chapteropening}

\section[die Bibliothek]{die Bibliothek — a library}

\label{lesson:GE-C98-bibliothek}



\begin{quote}
Before the new one: say the German for a village, then the German for the police.
\end{quote}

\subsection*{You'll want to know: die Bibliothek}
\textbf{die Bibliothek} — \textquotedblleft{}a library\textquotedblright{}.

\begin{cousinweb}
\textbf{Bibliothek} $\leftarrow$ Latin \emph{bibliothēca} $\leftarrow$ Greek \emph{biblioth\'{\={e}}kē}, built from \emph{biblíon}, \textquotedblleft{}book\textquotedblright{}, and \emph{th\'{\={e}}kē}, \textquotedblleft{}store\textquotedblright{} (EtymDB 2.1) — borrowed.
\end{cousinweb}

\begin{grammarlens}[title={What you've built}]
One more word for this chapter.
\end{grammarlens}

\subsection*{Your turn}
\begin{itemize}
\raggedright
  \item \emph{Say it:} \emph{die Bibliothek}
  \item \emph{Say it:} \emph{die Bibliothek}, once more
  \item \emph{Say it:} say \emph{die Bibliothek}
  \item \emph{Recall:} say the German for a village, then the German for the police, then say \emph{die Bibliothek} again
\end{itemize}

\begin{tcolorbox}[breakable,colback=teal!4,colframe=teal!35!black,title={Before you move on}]
What does die Bibliothek mean? (\textquotedblleft{}A library\textquotedblright{}.) Say it once more.
\end{tcolorbox}
\section[das Fahrrad]{das Fahrrad — a bicycle}

\label{lesson:GE-C98-fahrrad}



\begin{quote}
Before the new one: say the German for the police, then the German for a library.
\end{quote}

\subsection*{You'll want to know: das Fahrrad}
\textbf{das Fahrrad} — \textquotedblleft{}a bicycle\textquotedblright{}.

\begin{cousinweb}
\textbf{Fahrrad} is \emph{fahren}, \textquotedblleft{}to travel\textquotedblright{}, + the word for \textquotedblleft{}wheel\textquotedblright{}, a cousin of Latin \emph{rota}, \textquotedblleft{}wheel\textquotedblright{} (EtymDB 2.1 lists it beside English \emph{rotary}): a wheel you ride.
\end{cousinweb}

\begin{grammarlens}[title={What you've built}]
One more word for this chapter.
\end{grammarlens}

\subsection*{Your turn}
\begin{itemize}
\raggedright
  \item \emph{Say it:} \emph{das Fahrrad}
  \item \emph{Say it:} \emph{das Fahrrad}, once more
  \item \emph{Say it:} say \emph{das Fahrrad}
  \item \emph{Recall:} say the German for the police, then the German for a library, then say \emph{das Fahrrad} again
\end{itemize}

\begin{tcolorbox}[breakable,colback=teal!4,colframe=teal!35!black,title={Before you move on}]
What does das Fahrrad mean? (\textquotedblleft{}A bicycle\textquotedblright{}.) Say it once more.
\end{tcolorbox}
\section[die Haltestelle]{die Haltestelle — a bus or tram stop}

\label{lesson:GE-C98-haltestelle}



\begin{quote}
Before the new one: say the German for a library, then the German for a bicycle.
\end{quote}

\subsection*{You'll want to know: die Haltestelle}
\textbf{die Haltestelle} — \textquotedblleft{}a bus or tram stop\textquotedblright{}.

\begin{cousinweb}
\textbf{Haltestelle} is the verb \textquotedblleft{}to hold, to stop\textquotedblright{} (Old High German \emph{haltan} $\leftarrow$ Proto-Germanic \emph{*haldaną}, English \emph{hold}, EtymDB 2.1) + the word for \textquotedblleft{}place\textquotedblright{}: the stopping place.
\end{cousinweb}

\begin{grammarlens}[title={What you've built}]
One more word for this chapter.
\end{grammarlens}

\subsection*{Your turn}
\begin{itemize}
\raggedright
  \item \emph{Say it:} \emph{die Haltestelle}
  \item \emph{Say it:} \emph{die Haltestelle}, once more
  \item \emph{Say it:} say \emph{die Haltestelle}
  \item \emph{Recall:} say the German for a library, then the German for a bicycle, then say \emph{die Haltestelle} again
\end{itemize}

\begin{tcolorbox}[breakable,colback=teal!4,colframe=teal!35!black,title={Before you move on}]
What does die Haltestelle mean? (\textquotedblleft{}A bus or tram stop\textquotedblright{}.) Say it once more.
\end{tcolorbox}
\section[die Kleidung]{die Kleidung — clothes}

\label{lesson:GE-C98-kleidung}



\begin{quote}
Before the new one: say the German for a bicycle, then the German for a bus or tram stop.
\end{quote}

\subsection*{You'll want to know: die Kleidung}
\textbf{die Kleidung} — \textquotedblleft{}clothes\textquotedblright{}.

\begin{cousinweb}
\textbf{Kleidung} $\leftarrow$ Middle High German \emph{kleidunge}, from the verb \textquotedblleft{}to clothe\textquotedblright{} + \emph{-ung} (EtymDB 2.1). Its root is the word for \textquotedblleft{}dress\textquotedblright{}, Proto-Germanic \emph{*klaiþą}, \textquotedblleft{}cloth\textquotedblright{}, English \emph{cloth}.
\end{cousinweb}

\begin{grammarlens}[title={What you've built}]
One more word for this chapter.
\end{grammarlens}

\subsection*{Your turn}
\begin{itemize}
\raggedright
  \item \emph{Say it:} \emph{die Kleidung}
  \item \emph{Say it:} \emph{die Kleidung}, once more
  \item \emph{Say it:} say \emph{die Kleidung}
  \item \emph{Recall:} say the German for a bicycle, then the German for a bus or tram stop, then say \emph{die Kleidung} again
\end{itemize}

\begin{tcolorbox}[breakable,colback=teal!4,colframe=teal!35!black,title={Before you move on}]
What does die Kleidung mean? (\textquotedblleft{}Clothes\textquotedblright{}.) Say it once more.
\end{tcolorbox}
\section[das Hemd]{das Hemd — a shirt}

\label{lesson:GE-C98-hemd}



\begin{quote}
Before the new one: say the German for a bus or tram stop, then the German for clothes.
\end{quote}

\subsection*{You'll want to know: das Hemd}
\textbf{das Hemd} — \textquotedblleft{}a shirt\textquotedblright{}.

\begin{cousinweb}
\textbf{Hemd} $\leftarrow$ Middle High German \emph{hemde} $\leftarrow$ Old High German \emph{hemidi} $\leftarrow$ Proto-Germanic \emph{*hamiþiją}, \textquotedblleft{}clothing, shirt\textquotedblright{} — inherited; Dutch \emph{hemd} is the same word. (EtymDB 2.1)
\end{cousinweb}

\begin{grammarlens}[title={What you've built}]
One more word for this chapter.
\end{grammarlens}

\subsection*{Your turn}
\begin{itemize}
\raggedright
  \item \emph{Say it:} \emph{das Hemd}
  \item \emph{Say it:} \emph{das Hemd}, once more
  \item \emph{Say it:} say \emph{das Hemd}
  \item \emph{Recall:} say the German for a bus or tram stop, then the German for clothes, then say \emph{das Hemd} again
\end{itemize}

\begin{tcolorbox}[breakable,colback=teal!4,colframe=teal!35!black,title={Before you move on}]
What does das Hemd mean? (\textquotedblleft{}A shirt\textquotedblright{}.) Say it once more.
\end{tcolorbox}
